\section{Schlussfolgerungen und Empfehlungen}
Die durchgeführte Analyse der Teststrategien Testpyramide und Testing Honeycomb Modell hat gezeigt, dass beide Ansätze spezifische Vor- und Nachteile mit sich bringen. Während die Testpyramide durch ihre modulare Struktur eine höhere Testgeschwindigkeit und gezieltere Fehlerdiagnose ermöglicht, bietet das Testing Honeycomb Modell eine realistischere Abbildung der Systeminteraktionen, geht jedoch mit höheren Laufzeiten und einem komplexeren Debugging Prozess einher. 

Die Messungen zeigen, dass die Testpyramide insgesamt eine drei Mal schnellere Testdurchführung bietet, während das Testing Honeycomb Modell eine um fünfzehn Prozent höhere Systemabdeckung erreicht. Besonders auffällig ist die erhöhte Anzahl an Tests in der Testpyramide, da diese Strategie kleinere Codeeinheiten isoliert testet, während das Testing Honeycomb Modell vollständige End to End Abläufe prüft.  

Die Testpyramide erleichtert durch ihre schichtweise getesteten Komponenten die gezielte Fehlerdiagnose. Probleme können frühzeitig auf Methoden oder Serviceebene erkannt werden. Das Testing Honeycomb Modell erfordert oft eine vollständige Analyse des Anwendungsflusses, um Fehler zu lokalisieren. Dies macht Debugging deutlich aufwendiger. Besonders in der Mapper Schicht zeigte sich, dass Fehler in der Testpyramide schnell identifiziert wurden, während sie im Testing Honeycomb Modell nur durch tiefgehende Log Analysen auffindbar waren.  

Die Einführung von Spring Security hat in beiden Modellen zu einer erhöhten Anzahl an Testfällen und längeren Laufzeiten geführt. Während die Testpyramide durch ihren modularen Aufbau besser mit den zusätzlichen Sicherheitsprüfungen umgehen konnte, führten sie im Testing Honeycomb Modell zu einer deutlichen Laufzeitsteigerung. Die Ergebnisse zeigen, dass zusätzliche Sicherheitsmechanismen in End to End Tests erheblichen Overhead erzeugen und ihre Auswirkungen frühzeitig berücksichtigt werden sollten.  

Die Wahl der Datenbank hat ebenfalls eine signifikante Rolle gespielt. Tests mit H zwei liefen deutlich schneller als mit MySQL Testcontainern. Allerdings gab es Kompatibilitätsprobleme bei SQL Syntax und Fremdschlüsselvalidierungen. Besonders auffällig war die unterschiedliche Interpretation von Escape Zeichen und Fremdschlüssel Checks zwischen H zwei und MySQL, die zu unerwartetem Verhalten in Tests führten. MySQL Testcontainer bieten eine realitätsnähere Umgebung, benötigen jedoch signifikant längere Initialisierungszeiten, was insbesondere im Testing Honeycomb Modell problematisch war.  

Beim Testing Honeycomb Modell konnte durch Test Driven Development zunächst mit einem einzigen Test begonnen werden, ohne direkt Service oder Repository Klassen zu implementieren. Die Testpyramide erfordert eine schrittweise Implementierung der Schichten, erlaubt aber eine stabilere Codebasis und eine schnellere Testausführung. Nachzüglertests, also Tests die nachträglich hinzugefügt wurden, zeigten, dass die Testpyramide eine bessere Integration neuer Tests ermöglichte, während Testing Honeycomb Tests häufiger vollständig überarbeitet werden mussten.  

Die Ergebnisse legen nahe, dass eine hybride Teststrategie die Stärken beider Modelle kombinieren könnte. Eine mögliche Lösung wäre die Nutzung von Unit Tests und gezielten Integrationstests aus der Testpyramide für schnelle Entwicklungszyklen. Zusätzlich könnten End to End Tests aus dem Testing Honeycomb Modell für komplexe geschäftsprozessorientierte Abläufe integriert werden. Eine adaptive Teststrategie wäre ebenfalls eine Möglichkeit, um je nach Codeänderung automatisch die passende Testmethode auszuwählen.  

Zukünftige Arbeiten könnten sich darauf konzentrieren, wie eine automatisierte Auswahl der Teststrategie entwickelt werden kann, um je nach Systemänderung effizientere Testmechanismen einzusetzen. Eine weitere Optimierung der Datenbankkonfiguration könnte helfen, die Performance von MySQL Testcontainern zu verbessern und Escape Syntax Probleme zu minimieren. Ebenso könnte die detaillierte Untersuchung hybrider Testmodelle zeigen, wie Unit und Integrationstests mit End to End Tests kombiniert werden können, um eine bessere Balance zwischen Effizienz und Systemabdeckung zu erreichen.  

Die Wahl der Teststrategie sollte immer kontextabhängig erfolgen. Die Testpyramide eignet sich besonders für modularen Code und agile Entwicklungsprozesse, bei denen schnelle Rückmeldungen benötigt werden. Das Testing Honeycomb Modell bietet Vorteile für komplexe Systeminteraktionen, erfordert jedoch einen höheren Laufzeitaufwand. Eine Kombination beider Ansätze könnte langfristig eine effiziente und skalierbare Teststrategie ermöglichen, die sowohl die Testabdeckung als auch die Testgeschwindigkeit optimiert.  

Diese Arbeit liefert eine praxisnahe Evaluierung beider Testmodelle und zeigt auf, wie hybride Teststrategien zur Effizienzsteigerung und Fehlerprävention beitragen können.  


\subsection{Zusammenfassung der Ergebnisse je Anwendungsfall}
\subsection{Empfehlungen für den praktischen Einsatz der Teststrategien und Datenbanklösungen}
\subsection{Limitationen und zukünftige Arbeiten}  

Die durchgeführte Untersuchung beschränkt sich auf eine monolithische Architektur mit einer begrenzten Anzahl an Features. Zukünftige Arbeiten könnten untersuchen, wie sich die Teststrategien auf komplexere Geschäftsprozesse auswirken, insbesondere wenn zusätzliche Funktionalitäten integriert werden.  

Ein Aspekt, der in dieser Arbeit nicht behandelt wurde, ist die Anwendung des Testing Honeycomb Modells in einer Microservices Architektur. Eine Studie von Spotify \cite{
Spotify18} beschreibt, dass in Microservices Umgebungen End to End Tests eine zentrale Rolle spielen, da sie sicherstellen, dass die Kommunikation zwischen den unabhängigen Diensten zuverlässig funktioniert.  

Microservices sind ein Architekturstil, bei dem eine Anwendung in mehrere unabhängige Dienste unterteilt wird, die über Netzwerkschnittstellen interagieren. Diese Definition basiert auf den Arbeiten von Fowler und Lewis \cite{fowler2014} sowie Newman \cite{newman2015}, die Microservices als lose gekoppelte, eigenständig deploybare Einheiten beschreiben. Besonders in verteilten Systemen, in denen Services synchron über REST APIs oder asynchron über Messaging Systeme kommunizieren, ergeben sich besondere Herausforderungen bei der Testbarkeit. Während Unit und Integrationstests die isolierte Geschäftslogik eines einzelnen Services validieren, erfassen sie nicht potenzielle Fehler in der Kommunikation zwischen den Services. Das Testing Honeycomb Modell könnte in solchen Szenarien vorteilhaft sein, da es die gesamte Service Interaktion unter realistischen Bedingungen überprüft.  

Die Testergebnisse legen nahe, dass die Testpyramide für kleine bis mittelgroße monolithische Anwendungen eine effiziente Teststrategie darstellt. Sie ermöglicht eine gezielte Fehlerdiagnose und eine hohe Testabdeckung einzelner Komponenten. Mit zunehmender Größe eines Monolithen steigt jedoch die Anzahl der internen Abhängigkeiten, was die Testwartung erschwert. Zudem kann das Testing Honeycomb Modell bei sehr großen Monolithen zu unverhältnismäßig langen Testlaufzeiten führen. Dies deutet darauf hin, dass ab einer bestimmten Komplexitätsstufe eine Aufteilung in Microservices sinnvoll sein könnte, um die Testbarkeit zu verbessern.  

Ein vielversprechender Forschungsansatz könnte die Kombination der Testpyramide mit dem Testing Honeycomb Modell sein. Eine mögliche Strategie wäre die Nutzung der Testpyramide für isolierte Unit- und Integrationstests, um Fehler frühzeitig zu erkennen, während das Testing Honeycomb Modell gezielt für kritische Geschäftsprozesse eingesetzt wird. Anstatt das gesamte System in jedem Testdurchlauf vollständig zu validieren, könnten selektive Honeycomb-Tests durchgeführt werden, um Laufzeiten zu optimieren.  

Zukünftige Arbeiten könnten sich darauf konzentrieren, adaptive Teststrategien zu entwickeln, die je nach Art der Codeänderung automatisch zwischen granularen Tests der Testpyramide und umfassenden Systemtests des Testing Honeycomb Modells entscheiden. Besonders relevant wäre eine genauere Untersuchung, wie Testcontainer optimiert werden können, um Laufzeiten zu reduzieren und gleichzeitig eine hohe Testabdeckung sicherzustellen. Darüber hinaus könnte erforscht werden, inwiefern eine hybride Teststrategie für Microservices anwendbar ist, bei der die Testpyramide innerhalb einzelner Services genutzt wird und das Testing Honeycomb Modell für Service-übergreifende Tests zum Einsatz kommt.  
